\section{Experiments}
\label{sec:experiments}

\paragraph{Source and targets.}
Coswara cough-heavy and cough-shallow recordings form the source training set. COUGHVID is the largest public external target in our current audit. Tos COVID-19 is added as a second independent public external target; its 2021 official test split provides RT-PCR positive/negative labels and OGG cough clips collected by mobile phone \cite{buenosaires2021toscovid}. Virufy original and segmented clips are included only as tiny public stress targets, not as strong external validation cohorts. UK COVID-19 Vocal Audio and Cambridge COVID-19 Sounds / ComParE CCS remain high-priority larger follow-up targets, while DiCOVA is reserved as a Coswara-derived auxiliary protocol with overlap control.

\paragraph{Evaluation metrics.}
We report macro one-vs-rest AUROC, COVID one-vs-rest AUROC, macro AUPRC, ECE, Brier score, sensitivity at 95\% specificity, domain probe AUC, task probe AUC, checkpoint size, latency, and method-selection stability. Each primary COUGHVID result is averaged across seeds 7, 11, and 23.

\paragraph{Stress tests.}
COUGHVID is further analyzed with metadata-defined slices and bootstrap target subsets. Tos COVID-19 tests whether conclusions persist on a second independent public COVID cough target. Virufy segmented is evaluated both at the clip level and after subject-level probability aggregation. These stress tests ask whether the same KD method remains best when the target composition or evaluation unit changes.

\paragraph{Guard evaluation.}
We evaluate whether target-unlabeled signals can select the best KD strategy. Candidate signals include target confidence, entropy, predicted label distribution, and source-relative prediction shift. A guard is considered useful only if it selects a method that improves external metrics without increasing probe or calibration risk.
